\documentclass{jsarticle}
\usepackage{amssymb}
\usepackage{amsthm}
\usepackage{amsmath}
\usepackage{comment}
%定理環境 thmはあまり使わずpropをなるべく使う。
%主張は基本的に証明の中で使い、そのため番号を付けない。
%注意環境を付けてみた。
\newtheorem{thm}{定理}
\newtheorem{prop}[thm]{命題}
\newtheorem{cor}[thm]{系}
\newtheorem{lem}[thm]{補題}
\newtheorem{df}[thm]{定義}
\newtheorem{exam}[thm]{例}
\newtheorem{fact}[thm]{事実}
\newtheorem*{claim}{主張}
\newtheorem*{cau}{注意}
\renewcommand{\proofname}{{\bf 証明}}
%矢印たち。
%\toは\rightarrowと同じ。\getsは\leftarrowと同じ。同値は\iff
\newcommand{\To}{\Longrightarrow}
\newcommand{\Gets}{\LongLeftarrow}
%黒板太字
\newcommand{\nn}{\mathbb{N}}
\newcommand{\zz}{\mathbb{Z}}
\newcommand{\qq}{\mathbb{Q}}
\newcommand{\rr}{\mathbb{R}}
\newcommand{\cc}{\mathbb{C}}
%無限大
\newcommand{\mgn}{\infty}
%差集合は\setminusで統一する。等号の否定は\neq
%真部分集合は\subsetneqq　偏微分は\partial
%ディスプレイスタイル。\dfracでディスプレイ分数
\newcommand{\dpst}{\displaystyle}
\newcommand{\supp}{\text{{\rm supp}}}
%なんか演算２
\newcommand{\bd}{\text{{\rm Bdry}}}
\newcommand{\cl}{\text{{\rm CL}}}
\newcommand{\nai}{\text{{\rm Int}}}
%差集合は\setminusで統一する。等号の否定は\neq
%真部分集合は\subsetneqq　偏微分は\partial
%ディスプレイスタイル。\dfracでディスプレイ分数
%空集合は\Oを使う！！！！！！！！！！！！

\title{コンパクトハウスドルフ空間の正規性ほか}
\author{一色三良(concious77)}
\date{\today}
			\begin{document}
\maketitle

この文章ではコンパクトハウスドルフ空間の正規性と局所コンパクトハウスドルフ空間の完全正則性を証明する。その後コンパクトハウスドルフ空間に対する有限開被覆に従属する１の分割の存在を証明する。実は有限開被覆に従属する１の分割の存在は空間の正規性と同値なのであるが、それはまた別の話。以下では全ての空間にハウスドルフ性が仮定されているので正規という言葉と$T_4$という言葉、及び正則という言葉と$T_3$という言葉を特に区別しないことにする。

\begin{prop}
コンパクトハウスドルフ空間は正則である。
\end{prop}

\begin{proof}$X$をコンパクトハウスドルフ空間とし、$x\not\in F$となる点$x\in X$と閉集合$F$の任意の組を考える。$X$はハウスドルフ空間なので任意の$a\in F$に対して２つの開集合$O_a,P_a$が存在して
\[
x\in O_a, a\in P_a,O_a\cap P_a=\emptyset
\]
とを充たす。さて明らかに
\[
F\subset \bigcup_{a\in F}P_a
\]
であり。$X$のコンパクト性からその閉集合$F$もコンパクトになるので有限個の$a_1,a_2,\cdots a_n$が存在して
\[
F\subset \bigcup_{k=1}^{n}P_{a_k}
\]
となるこの右辺を$P$とおこう。つまり
\[
P=\bigcup_{k=1}^{n}P_{a_k} 
\]
これに対応して
\[
O=\bigcap_{k=1}^{n}O_{a_k}
\]
とすると$O$は開集合で$x\in O$であり、各$a_i$について
\[
O_{a_i}\cap P_{a_i}=\emptyset
\]
となり$O\subset O_{a_i}$なので任意の$a_i$について
\[
O\cap P_{a_i}=\emptyset
\]
更に$P=\bigcup_{k=1}^{n}P_{a_k} $だから
\[
O\cap P=\emptyset
\]
そして
$x\in O$で$F\subset P$なので$X$が正則であることが分かった。

\end{proof}

\begin{thm}
コンパクトハウスドルフ空間は正規である。
\end{thm}
\begin{proof}
$X$の互いに交わらない２つの閉集合を任意に固定して$A,B$とする。先の命題よりコンパクトハウスドルフ空間$X$は正則なので任意の$a\in A$毎に開集合$U_a,V_a$が存在して
\[
a\in U_a, B\subset V_a,U_a\cap V_a=\emptyset
\]
となる。さて明らかに
\[
A\subset \bigcup_{a\in A}U_a
\]
でありコンパクトハウスドルフ空間$X$の閉集合$A$もコンパクトなので有限個の$a_1,a_2\cdots a_n$が存在して
\[
A\subset \bigcup_{k=1}^{n}U_{a_k}
\]
となる。
\[
U=\bigcup_{k=1}^{n}U_{a_k}
\]
と定義しよう。
これに対応して
\[
V=\bigcap_{k=1}^{n}V_{a_k}
\]
とすると$V$は開集合であり、$B\subset V$であり各$a_i$で
\[
U_{a_i}\cap V_{a_i}=\emptyset
\]
で$V\subset V_{a_i}なので$任意の$a_i$について
\[
U_{a_i}\cap V=\emptyset
\]
これと$U=\bigcup_{k=1}^{n}U_{a_k}$から
\[
U\cap V=\emptyset
\]
を得る。
そして$A\subset U$と$B\subset V$なので$X$が正規であることが示せた。
\end{proof}

\begin{thm}
局所コンパクトハウスドルフ空間は完全正則である。
\end{thm}
\begin{proof}
局所コンパクトハウスドルフ空間$X$の互いに交わらない点と閉集合を$x,F$とし、$X$の１点コンパクト化を$Y$とし、無限遠点を$\mgn$とする。つまり$Y=X\cup\{\mgn\}$この時$G=F\cup \{\mgn\}$は$Y$の閉集合であり$x\not\in G$である。$Y$は正規なのでウリゾーンの補題から$f(x)=0,f|G=0$となる連続写像
$f:Y\to \rr$が存在する。これを$X$に制限すればそれも連続であり、$x$と$F$を分離する関数であるから$X$が完全正則であることがわかる。
\end{proof}

\begin{df}[関数の台]
写像$f:X\to \rr$の台を集合$\{x\in X\ |\ f(x)\neq 0\}$の閉包として定義し、これを$\supp f$で表す。
\end{df}
\begin{cau}
閉包を取ってることに注意せよ。
\end{cau}

%%%%%%%%%%%%%%%%%%%%%%%%%%%%%%%%%%%%%%%%%%%%%%%%%%%%%%%%%%%%%%%%%%%%%%
\begin{thm}[被覆の縮小の存在]\label{cover}
コンパクトハウスドルフ空間$X$の有限開被覆$\{U_k\}_{k=1,2,\cdots n}$
について有限開被覆$\{V_k\}_{k=1,2,\cdots n}$が存在して
\[
\cl(V_k)\subset U_k
\]
を充たす。
\end{thm}

\begin{proof}
帰納的に$V_k$を構成する。\\ 
まず$V_1$を構成する。$F_1=X\setminus \bigcup_{k\neq1}U_k$と置くとこれは閉集合で$\{U_k\}_{k=1,2,\cdots n}$が被覆を成すことから
\[
F_1\subset U_1
\]
となる。ここで$X$の正規性を用いると
\[
F_1\subset V\subset \cl(V)\subset U_1
\]
となる開集合$V$が存在する。このような$V$の一つを$V_1$とおくと
\[
\cl(V_1)\subset U_1
\]
を明らかに充たし
更に$\{V_1,U_2,U_3,\cdots U_n\}$が$X$の被覆を成している。\\ 
$V_k$まで
\[
\cl(V_k)\subset U_k
\]
を充たし$\{V_1,V_2,\cdots V_k,U_{k+1},U_{k+2}\cdots U_n\}$が被覆を成すような$V_k$が構成されたとしよう。この時
\[
F_{k+1}=X\setminus \left( \bigcup_{l=1}^{k}V_l\cup\bigcup_{l=k+2}^{n}U_l \right)
\]
とする（$U_{k+1}$は和集合されてる中にはないことに注意せよ）このとき$F_{k+1}$は閉集合で
$\{V_1,V_2,\cdots V_k,U_{k+1},U_{k+2}\cdots U_n\}$が被覆を成してるということから
\[
F_{k+1}\subset U_{k+1}
\]
となる。$X$の正規性を用いると
\[
F_{k+1}\subset V\subset \cl(V) \subset U_{k+1}
\]
となる。開集合$V$が存在するがこれの一つを$V_{k+1}$とすると
\[
\cl(V_{k+1}) \subset U_{k+1}
\]
を充たし$\{V_1,V_2,\cdots V_k,V_{k+1},U_{k+2}\cdots U_n\}$は被覆を成す。この構成を$n$まで行えば求める$\{V_k\}$が得られる。
\end{proof}
\begin{cau}
証明の中で$X$の正規性以外の性質を用いてないことに注意せよ。つまり、この定理は一般に正規空間でも成り立つのだが、それはまた別のお話。
\end{cau}
%%%%%%%%%%%%%%%%%%%%%%%%%%%%%%%%%%%%%%%%%%%%%%%%%%%%%%%%%%%%%%%%%%%%%%%%%%%%

\begin{thm}[被覆の縮小の存在2]\label{cover2}
コンパクトハウスドルフ空間$X$の有限開被覆$\{U_k\}_{k=1,2,\cdots n}$
について有限閉被覆$\{L_k\}_{k=1,2,\cdots n}$が存在して
\[
L_k\subset U_k
\]
を充たす。
\end{thm}
\begin{proof}
$\{U_k\}_{k=1,2,\cdots n}$に対して定理\ref{cover}を用いて
\[
\cl(V_k)\subset U_k
\]
を充たす有限開被覆$\{V_k\}_{k=1,2,\cdots n}$が存在する。これを用いて
\[
L_k=\cl(V_k)
\]
とすれば良い。
\end{proof}
\begin{cau}定理\ref{cover2}は有限閉（closed）被覆の存在を言ってることに注意せよ。開（open）被覆ではない。
\end{cau}

\begin{thm}[コンパクトハウスドルフ空間上の単位の分割]\label{pat}
コンパクトハウスドルフ空間$X$の有限開被覆$\{U_k\}_{k=1,2,\cdots n}$
について実連続関数の族$\{f_k\}_{k=1,2,\cdots n}$が存在して以下を充たす。
\begin{eqnarray*}
\supp f_k\subset U_k\\ 
0\le f_k \le1\\ 
\sum_{k=1}^{n}f_k\equiv 1
\end{eqnarray*}
\end{thm}
\begin{proof}
定理\ref{cover2}から$\{U_k\}_{k=1,2,\cdots n}$に対して
\[
L_k\subset U_k
\]
を充たす有限閉被覆$\{L_k\}_{k=1,2,\cdots n}$が存在する。ここでウリゾーンの補題より各$k$毎に連続関数$g_k:X\to [0,1]$が存在して
\[
g_k|L_k=1,\ g_k|X\setminus U_k=0
\]
を充たす。つまり
\[
\supp (g)\subset U_k
\]
である。
$\{L_k\}_{k=1,2,\cdots n}$が被覆であることから$\dpst \sum_{k=1}^ng_k$は常に正である。そこで
\[
f_k=\frac{g_k}{\sum_{k=1}^ng_k}
\]と置けばこれは連続で明らかに
\begin{eqnarray*}
\supp f_k\subset U_k\\ 
0\le f_k \le1\\ 
\sum_{k=1}^{n}f_k\equiv 1
\end{eqnarray*}
を充たす。
\end{proof}


\begin{cor}[局所コンパクトハウスドルフ空間のコンパクト集合上の単位の分割]\label{pat}
局所コンパクトハウスドルフ空間$X$のコンパクト集合$K$と有限個の開集合の族
$\{U_k\}_{k=1,2,\cdots n}$が
\[
K\subset \bigcup_{k=1}^nU_k
\]を充たしてるとする。このとき実連続関数の族$\{f_k\}_{k=1,2,\cdots n}$が存在して以下を充たす。
\begin{eqnarray*}
&\supp f_k\subset U_k\\ 
&0\le f_k \le1\\ 
&\forall x\in K\ \sum_{k=1}^{n}f_k(x)\equiv 1
\end{eqnarray*}


\end{cor}
\begin{proof}
$X$の１点コンパクト化を$Y$と置くこの時$\{U_1,U_2,\cdots U_n,Y\setminus K\}$は$Y$の被覆となっている。都合のために$U_{n+1}=Y\setminus K$と置こう。$\{U_k\}_{k=1,2,\cdots n,n+1}$は$Y$の被覆を成してるので先の定理\ref{pat}より
\begin{eqnarray*}
&\supp f_k\subset U_k\\ 
&0\le f_k \le1\\ 
&\forall x\in K\ \sum_{k=1}^{n+1}f_k(x)\equiv 1
\end{eqnarray*}
を充たす実連続関数の族$\{f_k\}_{k=1,2,\cdots n,n+1}$が存在する。ここで
\[
\supp f_{n+1}\subset U_{n+1}=Y\setminus K
\]
より
\[
f_{n+1}|K\equiv0
\]
よって$\{f_k\}_{k=1,2,\cdots n,n+1}$から$f_{n+1}$を除いた$\{f_k\}_{k=1,2,\cdots n}$は明らかに
\begin{eqnarray*}
&\supp f_k\subset U_k\\ 
&0\le f_k \le1\\ 
&\forall x\in K\ \sum_{k=1}^{n}f_k(x)\equiv 1
\end{eqnarray*}
を充たす。
\end{proof}

\begin{df}
位相空間$X$の部分集合$A$が相対コンパクトであるとはその閉包$\cl(A)$がコンパクトである時に言う
\end{df}

\begin{lem}
局所コンパクトハウスドルフ空間$X$とそのコンパクト集合$K$と開集合$G$について$K\subset G$が充たされるならば、相対コンパクトな開集合$V$が存在して
\[
K\subset V\subset \cl(V) \subset G
\]
を充たす。つまり局所コンパクトハウスドルフ空間のコンパクト集合は相対コンパクトな基本近傍系を持つということである。
\end{lem}
\begin{proof}
$X$は局所コンパクトハウスドルフなので$K$も各点$a\in K$について相対コンパクトな開集合$V_a$が存在して
\[
V_a\subset \cl(V_a)\subset G
\]
を充たす。また明らかに
\[
K\subset \bigcup_{a\in K}V_a
\]
なので$K$のコンパクト性から有限個の$a_1,a_2,\dots ,a_k$が存在して
\[
K\subset \bigcup_{i=1}^{k}V_{a_i}
\]
となる。
\[
\cl(\bigcup_{i=1}^{k}V_{a_i})=\bigcup_{i=1}^{k}\cl(V_{a_i})
\]
なので$\dpst V=\bigcup_{i=1}^{k}V_{a_i}$も相対コンパクトで
\[
V\subset \cl(V)\subset G
\]
なので定理の証明は終わる。
\end{proof}

\begin{thm}[局所コンパクトハウスドルフ空間のコンパクト集合上のコンパクトサポートな単位の分割]
局所コンパクトハウスドルフ空間$X$のコンパクト集合$K$と有限個の開集合の族
$\{U_k\}_{k=1,2,\cdots n}$が
\[
K\subset \bigcup_{k=1}^nU_k
\]を充たしてるとする。このとき実連続関数の族$\{f_k\}_{k=1,2,\cdots n}$が存在して以下を充たす。
\begin{eqnarray*}
&\text{各$\supp f_k$はコンパクト}\\
&\supp f_k\subset U_k\\ 
&0\le f_k \le1\\ 
&\forall x\in K\ \sum_{k=1}^{n}f_k(x)\equiv 1
\end{eqnarray*}


\end{thm}
\begin{proof}
$K$と$\dpst \bigcup_{k=1}^nU_k$に対して先の補題を用いれば相対コンパクトな開集合$V$が存在して
\[
K\subset V\subset \cl(V)\subset \bigcup_{k=1}^nU_k
\]
を充たす。\par
さて$Ok=U_k\cap V$と置くと、各$O_k$は相対コンパクトであり、
\[
K\subset \bigcup_{k=1}^nO_k
\]
を充たす。$K$と$\{O_k\}_{k=1,2,\cdots n}$に対して系\ref{pat}を用いると、
実連続関数の族$\{f_k\}_{k=1,2,\cdots n}$が存在して以下を充たす。
\begin{eqnarray*}
&\supp f_k\subset O_k\\ 
&0\le f_k \le1\\ 
&\forall x\in K\ \sum_{k=1}^{n}f_k(x)\equiv 1
\end{eqnarray*}
ところで$O_k$が相対コンパクトであることから各$\supp f_k$はコンパクトであり、$O_k\subset U_k$なので定理は示された。
\end{proof}






	
\begin{thebibliography}{9}
\bibitem{1}森田紀一,位相空間論,岩波全書331,岩波書店,1981
\bibitem{2}W.Rudin,Real and Complex Analysis second edition,McGraw-Hill, 1974
\end{thebibliography}




			\end{document}



\begin{comment}

[字体の変更]

{\rm hogehoge}と使う。

\rm ローマン
\it イタリック
\sl 斜体

[supp]

\newcommand{\supp}{\text{{\rm supp}}}

[中央揃えの改行]

	\begin{eqnarray*}
		hoge1&=&hogehoge
		hoge2&=&hogehofe

	\end{eqnarray*}

&&の部分で揃えられる。






[場合分け]

	\begin{cases}
		hoge1 & \text{状況１}\\
		hoge2 & \text{状況２}\\
		・
		・
		・
	\end{cases}



\end{comment}





